\chapter{Desenvolvimento}

Este capítulo detalha o funcionamento de algumas funcionalidades
desenvolvidas durante o projeto. Apenas algumas funcionalidades mais
chaves são abordadas neste capítulo, como o registro de humor, vinculo
de um gatilho, indicadores, e o registro de atividades.

\section{Processo de Login}

Logo ao abrir o aplicativo, o usuário é redirecionado para a tela de
login. A tela em si é relativamente simples, e contem os campos de
email e senha. Existe um elemento proposital aqui, uma forma de gerar
um certo nível de anonicidade, onde o usuário não precisa se identificar
com um nome ou número de telefone, ou mesmo ativar através de uma conta
de outra plataforma, como por exemplo, Google ou Facebook.

Em especial por dois motivos, (1) necessitaria uma conexão OAuth, de
forma que as plataformas de terceiro poderiam rastrear dados de qual
usuário tem acesso a qual aplicativo. Isso é importante porque existe,
e é fartamente documentado, sobre os mecanismos usados por tais plataformas
para rastrear e monitorar o acesso aos dados dos usuários, de forma a
vender e monetizar os dados dos usuários.

% TODO: Write citations, link to news articles about it

Por fim, ainda retomando a proposta de oferecer a possibilidade de
\textit{self-service} para o usuário, o usuário que decidisse
ter uma instância do aplicativo própria para si, não precisaria de uma
conta em uma plataforma de terceiro, como Google ou Facebook e tampouco
precisaria registrar-se nessas plataformas.

\begin{figure}[H]
\centering
\caption{Tela de login}
\label{fig:screenshot-login}
\includegraphics[width=0.70\textwidth,height=0.60\textheight,keepaspectratio]{screenshot-login.jpeg}
\source{O autor, 2026}
\end{figure}

Por fim, destaca-se aqui que, conforme mencionado anteriormente, na
modelagem de casos de uso \ref{fig:use-cases}, ao realizar o login,
no primeiro acesso é enviado um email com um código de verificação,
seguindo a logistica de \textit{magic links}, onde adiciona uma camada
extra de segurança, garantindo que apenas o usuário que solicitou o login
possa acessar a conta.

\section{Página de Consulta de Registros}

A tela inicial é a tela nomeada apenas como ``Novo''. A proposta aqui
é ser a tela de cadastro rápido de um humor, e também é a tela que
visualmente mais destoa das demais. O usuário é apresentado com três
seções distintas:

\begin{itemize}
    \item Cadastro rápido de um humor
    \item Entre ontem e hoje
    \item Registros recentes
\end{itemize}

Sendo respectivamente, uma seção onde o usuário é convidado a
``escolher um humor'', e iniciar o processo de registro de uma nova entrada
de humor, a seção onde é computado os indicadores do nível de energia
acumulado entre o dia de ontem e hoje, além da informação da média de
sono do dia atual e comparada com o dia anterior. Por fim, também é
listado os últimos dez registros de humor que o usuário cadastrou.

\begin{figure}[H]
\centering
\caption{Tela de Consulta de Registros}
\label{fig:screenshot-consulta-registros}
\includegraphics[width=0.70\textwidth,height=0.60\textheight,keepaspectratio]{screenshot-initial-page.jpeg}
\source{O autor, 2026}
\end{figure}

Aqui já é introduzido para o usuário um padrão de organização e hierarquia:
todos elementos centrais são organizados em \textit{cards}, que são
exibidos em no máximo duas colunas. Também é introduzido, caso o usuário
já tenha registros de humor cadastrados, como será resumido esses registros
cadastrados pelo usuário em todas as seções do aplicativo: geralmente,
tendo um ícone, um título e alguma informação adicional destacada.

\section{Criação de Registros de Humor}

Nesta seção, o usuário pode criar novos registros de humor, escolhendo
entre uma lista de opções pré-definidas. O humor selecionado na tela
anterior é mostrado aqui, e o usuário é convidado a julgar como ele se
sente referente a uma escala de 1 a 10, sobre os níveis de energia,
estresse e ansiedade, respectivamente.

\begin{figure}[H]
\centering
\caption{Criação de Registros de Humor}
\label{fig:screenshot-consulta-humor}
\includegraphics[width=0.70\textwidth,height=0.60\textheight,keepaspectratio]{screenshot-create-mood.jpeg}
\source{O autor, 2026}
\end{figure}

Naturalmente, é uma questão dificil de responder, porém é importante
que seja um momento reflexivo e tampouco, como descreve o aplicativo,
não necessariamente sendo preciso.

O usuário pode adicionar uma anotação sobre aquele registro em
específico, que igualmente é um padrão que se repete/apresenta durante
outras telas.

Por fim, o usuário pode escolher "o que compõe aquele registro", que
ou seja, escolher emoções e estados/sentimentos que ele sente naquele
momento. Aqui uma listagem mista desses elementos é apresentado, e que
cujo ele pode selecionar os que mais se aplica a ele, e julgar um nível
de intensidade para cada um. Aqui a escala é menor e mais fácil, visando
facilitar a escolha e julgamento do usuário, entre ``Leve'', ``Moderado'',
e ``Intenso'', conforme a figura \ref{fig:screenshot-add-components}.

\begin{figure}[H]
\centering
\caption{Adicionando sentimentos/emoções}
\label{fig:screenshot-add-components}
\includegraphics[width=0.70\textwidth,height=0.60\textheight,keepaspectratio]{screenshot-add-components.jpeg}
\source{O autor, 2026}
\end{figure}

Ao concluir esse processo, e salvar, é apresentado para o usuário
uma tela de confirmação, onde ele pode opcionalmente vincular
esse registro a um gatilho em específico, que pode ter ocorrido
antes de ele salvar o registro, ou opcionalmente, criar um novo
gatilho. Aqui o projeto novamente, se ancora no EMA, recheando o registro
de humor com mais contexto sobre o que aconteceu.

\begin{figure}[H]
\centering
\caption{Vinculando um humor com um gatilho}
\label{fig:screenshot-link-mood-trigger}
\includegraphics[width=0.70\textwidth,height=0.60\textheight,keepaspectratio]{screenshot-link-mood-with-trigger.jpeg}
\source{O autor, 2026}
\end{figure}

Posteriormente, o usuário pode delimitar o impacto que aquele gatilho
pode ter exercído naquele registro de humor em especial. Aqui, é,
novamente o processo de chute de impacto, nada com muita precisão.

\begin{figure}[H]
\centering
\caption{Avaliando o impacto de um gatilho no humor}
\label{fig:screenshot-assess-trigger-mood-impact}
\includegraphics[width=0.70\textwidth,height=0.60\textheight,keepaspectratio]{screenshot-link-mood-assess-impact.jpeg}
\source{O autor, 2026}
\end{figure}

Por fim, vale mencionar que após esse processo, internamente, é enfileirado
um \textit{background job} que processa alguns indicadores: recalcular
o nível de energia daquele dia, apresentado na figura \ref{fig:screenshot-consulta-registros},
recalcular o indicador de correlação entre humor e energia, e agendar
o \textit{background job} que recalculará o nível de energia daquele
período (os últimos sete dias), apresentado na figura \ref{fig:screenshot-consulta-registros}.

\section{Indicadores}

Nessa tela é apresentado os indicadores construídos a partir dos
registros de humor, sono, gatilho. Além disso, embora em estado experimental,
também é apresentado um ``resumo semanal'' que condensa o mais central
de todos os indicadores, apresentando o nível de humor, média de sono
total do período.

\begin{figure}[H]
\centering
\caption{Tela de indicadores}
\label{fig:screenshot-insights}
\includegraphics[width=0.70\textwidth,height=0.60\textheight,keepaspectratio]{screenshot-insights.jpeg}
\source{O autor, 2026}
\end{figure}

\section{Histórico}

Nessa tela é apresentado o histórico de registros de humor, sono, gatilho,
e ações de cuidado. Como os dados são diferentes, é consultado os
\textit{endpoints} respectivos de cada dado, paralelamente, e então é
processado no aplicativo.

\begin{figure}[H]
\centering
\caption{Tela de Histórico}
\label{fig:screenshot-history}
\includegraphics[width=0.70\textwidth,height=0.60\textheight,keepaspectratio]{screenshot-history.jpeg}
\source{O autor, 2026}
\end{figure}

Para fazer isso, é utilizado um pequeno serviço que atua com o padrão
de \textit{mappers} para transformar os dados de cada \textit{endpoint}
em um formato comum, que é utilizado pelo restante da tela.

\begin{lstlisting}[caption={Formato comum da tela de histórico}, label={lst:history-card}]
export type HistoryCard = {
  id: string; // "<category>-<id>" to avoid collisions
  category: HistoryCategory;
  subtype?: CareActionSubtype; // only set when category === "care_action"
  timestamp: string; // ISO string — used for sorting and display
  title: string;
  summary: string;
  badge?: HistoryBadge;
  raw: MoodEntry | SleepRecord | Trigger | CareAction;
};
\end{lstlisting}

Conforme pode-se ver no código \ref{lst:history-card}, o formato comum
utilizado na tela de histórico, mapeia em essencialmente quatro propriedades
do dado original, além de uma propriedade adicional para identificar a categoria.

A propriedade \texttt{raw} é utilizada para armazenar o dado original,
permitindo que seja recuperado posteriormente quando necessário. A
propriedade \texttt{badge} é utilizada para armazenar qual dado é mostrado
como um ``indicador'' mais importante para aquele dado em específico,
por exemplo, que um determinado registro de sono foi ``Insuficiente'', se for
menor que 7 horas, mostrar se existem vínculos com aquele registro em
especial, ou mostrar qual nível de um determinado registro de humor
é considerado importante prestar atenção, como ``Ansiedade alta'',
``Energia alta'', etc. O último é importante também saber, por que uma
pessoa com TAB, pode estar com a energia alta e significar que ela está
em um período de mania, o que é útil para se detectar alguma alteração
ou mudança de ciclos.

Vale notar também que o badge é similar ao \textit{badge} que aparece na
tela de consulta de registros, apresentado na figura \ref{fig:screenshot-consulta-registros}.

\FloatBarrier
